\section{Methods}\label{sec:methods}
This toolbox operates on either skeletal animation sequences with an explicit hierarchy or point light display (PLD) trajectories represented as joint-position sequences. After import, both data types are converted into a common parent-space representation so that the same decomposition and editing operators can be applied.

\subsection{Animation data import and pre-processing}\label{subsec:methods-import}
For skeletal animation data, editable bones are inferred automatically from the skeleton tree by retaining bones that possess both a valid local quaternion track and a non-zero distal child offset. Let $\mathbf{v}_{i,0}\in\mathbb{R}^3$ denote bone $i$'s rest vector from the parent joint to the distal joint in parent coordinates. Given the local quaternion track $q_i(t)$, the parent-space distal-joint position is
\begin{equation}
\mathbf{r}_i(t)=\mathbf{R}\!\left(q_i(t)\right)\mathbf{v}_{i,0},
\end{equation}
where $\mathbf{R}(q_i(t))$ is the rotation matrix associated with $q_i(t)$. The corresponding bone length is $L_i=\|\mathbf{v}_{i,0}\|$. For visualization, an PLD preview can also be generated by sampling distal-joint world positions at every animation key time.

For PLD data, a CSV file with marker name, frame index, and 3D coordinates is parsed. Coordinates are converted from centimeters to meters and remapped to the internal Three.js axis convention as
\begin{equation}
\mathbf{g}_i(t)=\bigl[x_i^{\mathrm{csv}}(t),\,z_i^{\mathrm{csv}}(t),\,y_i^{\mathrm{csv}}(t)\bigr]^\top /100.
\end{equation}
Each marker is then assigned a parent marker $p(i)$ through the binding interface. The toolbox constructs a layer-wise update order $\{\mathcal{U}_k\}_{k=1}^{K}$ by breadth-first traversal of the parent graph. For root markers ($p(i)=0$), local and global coordinates are identical, i.e.
\begin{equation}
\boldsymbol{\ell}_i(t)=\mathbf{g}_i(t).
\end{equation}
For non-root markers,
\begin{equation}
\boldsymbol{\ell}_i(t)=\mathbf{g}_i(t)-\mathbf{g}_{p(i)}(t), \qquad
L_i(t)=\|\boldsymbol{\ell}_i(t)\|.
\end{equation}

For markers in the first two update layers, the parent-space basis is set to the identity basis and therefore $\mathbf{r}_i(t)=\boldsymbol{\ell}_i(t)$. For deeper layers, a moving orthonormal basis $\mathbf{B}_i(t)=\left[\hat{\mathbf{x}}_i(t)\ \hat{\mathbf{y}}_i(t)\ \hat{\mathbf{z}}_i(t)\right]$ is generated from the parent local vector. Its first axis is aligned with the parent vector,
\begin{equation}
\hat{\mathbf{x}}_i(t)=\frac{\boldsymbol{\ell}_{p(i)}(t)}{\|\boldsymbol{\ell}_{p(i)}(t)\|}.
\end{equation}
At the first frame, $\hat{\mathbf{y}}_i(1)$ and $\hat{\mathbf{z}}_i(1)$ are obtained by orthogonalizing an auxiliary vector against $\hat{\mathbf{x}}_i(1)$; at subsequent frames they are updated by the Rodrigues rotation that aligns the reference $\hat{\mathbf{x}}$-axis with the current one. The marker is then expressed in parent-space coordinates as
\begin{equation}
\mathbf{r}_i(t)=\mathbf{B}_i^\top(t)\boldsymbol{\ell}_i(t).
\end{equation}

If temporal extension is requested, both skeletal and PLD clips are repeated seamlessly by concatenating repeated tracks while removing the duplicated boundary sample between adjacent copies.

\subsection{Motion decomposition and parameter definition}\label{subsec:methods-decomposition}
For each bone or marker $i$, the parent-space trajectory $\mathbf{r}_i(t)$ over $T$ frames is arranged as a matrix $\mathbf{X}_i\in\mathbb{R}^{T\times 3}$. The toolbox applies uncentered principal component analysis (PCA) to $\mathbf{X}_i$, yielding eigenvectors $\mathbf{e}_{i,1}$, $\mathbf{e}_{i,2}$, $\mathbf{e}_{i,3}$ and framewise scores $s_{i,k}(t)$. The two components with the largest explained variance define the main plane (MP), and the remaining component defines the orthogonal $V_3$ axis:
\begin{equation}
c_{i,1}(t)=s_{i,v_1}(t), \qquad
c_{i,2}(t)=s_{i,v_2}(t), \qquad
c_{i,3}(t)=s_{i,v_3}(t).
\end{equation}
The MP angle and its in-plane projection radius are
\begin{equation}
\theta_i(t)=\operatorname{unwrap}\!\left(\operatorname{atan2}\!\bigl(c_{i,2}(t),c_{i,1}(t)\bigr)\right),
\qquad
\rho_i(t)=\sqrt{c_{i,1}^2(t)+c_{i,2}^2(t)},
\end{equation}
where angle unwrapping removes $2\pi$ discontinuities between adjacent frames. The toolbox stores the mean and zero-mean MP components as
\begin{equation}
\bar{\theta}_i=\frac{1}{T}\sum_{t=1}^{T}\theta_i(t), \qquad
\tilde{\theta}_i(t)=\theta_i(t)-\bar{\theta}_i.
\end{equation}
The out-of-plane component is represented as
\begin{equation}
z_i(t)=c_{i,3}(t), \qquad
\bar{z}_i=\frac{1}{T}\sum_{t=1}^{T} z_i(t), \qquad
\tilde{z}_i(t)=z_i(t)-\bar{z}_i.
\end{equation}

The periodic part of the motion is analyzed in the frequency domain by applying a discrete Fourier transform to $\tilde{\theta}_i(t)$. The implementation uses a sampling frequency of $30$~Hz for skeletal animation and $100$~Hz for imported PLD data. Before transformation, the signal is zero-padded to the next power of two. The editable parameters associated with each unit are MP amplitude scale, MP mean addition, optional frequency-domain filtering (none, low-pass, high-pass, band-pass, or band-stop, with user-defined cutoff frequencies, transition bandwidth, and stop-band residual amplitude ratio), $V_3$ amplitude scale, $V_3$ mean addition, and global phase shift. The same parameterization can be applied either to a single target or to a user-selected set of targets through the general manipulation interface.

\subsection{Motion manipulation and reconstruction}\label{subsec:methods-reconstruction}
Let $\mathcal{F}\{\tilde{\theta}_i\}$ denote the Fourier spectrum of the MP-amplitude signal. MP amplitude editing first rescales the complex spectrum by a factor $\alpha_i$, after which an optional filter gain $H_i(f)$ is applied:
\begin{equation}
\hat{\tilde{\theta}}_i'(f)=\alpha_i\,H_i(f)\,\hat{\tilde{\theta}}_i(f).
\end{equation}
Within the transition bands, $H_i(f)$ is interpolated with a raised-cosine taper so that the gain changes smoothly between pass and stop regions. The edited MP amplitude is recovered by inverse Fourier transform,
\begin{equation}
\tilde{\theta}_i'(t)=\mathcal{F}^{-1}\!\left\{\hat{\tilde{\theta}}_i'(f)\right\},
\end{equation}
and the edited MP angle becomes
\begin{equation}
\theta_i'(t)=\tilde{\theta}_i'(t)+\bar{\theta}_i', \qquad
\bar{\theta}_i'=\bar{\theta}_i+\Delta\theta_i.
\end{equation}

The $V_3$ component is manipulated by shifting its mean and scaling the reconstructed signal:
\begin{equation}
\bar{z}_i'=\bar{z}_i+\beta_i\bar{L}_i, \qquad
z_i'(t)=\gamma_i\bigl(\tilde{z}_i(t)+\bar{z}_i'\bigr),
\end{equation}
where $\bar{L}_i=L_i$ for skeletal bones and $\bar{L}_i=T^{-1}\sum_{t=1}^{T}L_i(t)$ for PLD markers. To preserve segment length, the edited $V_3$ displacement is clipped by the instantaneous segment length $L_i(t)$, with $L_i(t)\equiv L_i$ for skeletal bones, and the in-plane radius is recomputed as
\begin{equation}
\rho_i'(t)=\sqrt{\max\!\left(0,\,L_i^2(t)-z_i'^2(t)\right)}.
\end{equation}
The edited parent-space trajectory is then reconstructed in the original PCA basis:
\begin{equation}
\mathbf{r}_i'(t)=
\rho_i'(t)\cos\theta_i'(t)\,\mathbf{e}_{i,v_1}+
\rho_i'(t)\sin\theta_i'(t)\,\mathbf{e}_{i,v_2}+
z_i'(t)\,\mathbf{e}_{i,v_3}.
\end{equation}

For skeletal animation, the edited parent-space direction
\begin{equation}
\hat{\mathbf{d}}_i'(t)=\frac{\mathbf{r}_i'(t)}{\|\mathbf{r}_i'(t)\|}
\end{equation}
is converted back to a local quaternion track. Let $\hat{\mathbf{d}}_i(t)$ be the original parent-space direction and $q_i(t)$ the original local quaternion. The toolbox computes the minimal rotation $q_{\mathrm{dir},i}(t)$ that maps $\hat{\mathbf{d}}_i(t)$ to $\hat{\mathbf{d}}_i'(t)$, and updates the quaternion as
\begin{equation}
q_i'(t)=q_{\mathrm{dir},i}(t)\,q_i(t).
\end{equation}
This preserves the original axial twist while enforcing the edited bone direction. The resulting quaternion sequence directly overwrites the corresponding track in the current clip.

For PLD data, reconstruction proceeds hierarchically according to the update order. For markers in the first two update layers, the edited local trajectory is set directly as
\begin{equation}
\boldsymbol{\ell}_i'(t)=\mathbf{r}_i'(t).
\end{equation}
For deeper layers, a new moving basis $\mathbf{B}_i'(t)$ is recomputed from the edited parent local trajectory using the same Rodrigues update rule, and the edited local vector becomes
\begin{equation}
\boldsymbol{\ell}_i'(t)=\mathbf{B}_i'(t)\mathbf{r}_i'(t).
\end{equation}
Global marker trajectories are then obtained recursively:
\begin{equation}
\mathbf{g}_i'(t)=
\begin{cases}
\boldsymbol{\ell}_i'(t), & p(i)=0,\\
\mathbf{g}_{p(i)}'(t)+\boldsymbol{\ell}_i'(t), & p(i)\neq 0.
\end{cases}
\end{equation}
Finally, a phase-shift parameter $\phi_i\in[0,1]$ circularly shifts the edited sequence by $\lfloor \phi_iT\rfloor$ frames. In the skeletal pipeline the phase shift is applied to the reconstructed quaternion track, whereas in the PLD pipeline it is applied to the reconstructed parent-space trajectory before hierarchical accumulation.

\subsection{Animation data export}\label{subsec:methods-export}
After editing, skeletal animation is exported as a binary \texttt{.glb} file containing the original model and the current edited animation clip. PLD data are exported as a \texttt{.csv} file by writing every current position track as marker-wise frame coordinates. During CSV export, the toolbox applies the inverse unit conversion and axis remapping so that internal coordinates are converted back to the external PLD format.

The toolbox can additionally render stimulus videos in \texttt{.webm} format from Pane A, Pane B, or both panes. For video export, the user specifies one or multiple camera configurations $(\mathrm{yaw},\mathrm{pitch},\mathrm{distance},\mathrm{moveX},\mathrm{moveY})$; the toolbox then renders the animation at a fixed frame rate and resolution, and multiple outputs are bundled automatically into a \texttt{.zip} archive.
